\preface[Acknowledgements]
\label{sec:acknowledgements}
\addcontentsline{toc}{chapter}{Acknowledgements}

This book owes its existence, in part, to Charles Favre. Over the past two years, he has repeatedly commented on my earlier book \cite{XiaSampleBook} in terms I found insulting. I have adopted none of his suggestions. These exchanges nevertheless made me aware of a need: some mathematicians working in other fields wish to understand recent developments in pluripotential theory, but lack a satisfactory textbook through which to learn its foundations.

The classical literature presents a further difficulty. It contains numerous errors and gaps, many of which are known only to a small circle of experts. For someone learning the subject, these are substantial obstacles. A reader may spend considerable time trying to prove an incorrect statement or reconstruct an argument whose missing steps are familiar only to specialists. I therefore decided to write a book that systematically corrects these errors, fills the gaps, and develops the foundational material in a coherent and self-contained manner.

The book also resolves some questions left open in classical treatments. The most important concerns the definition of a plurisubharmonic function, discussed in \cref{chap:q02}. We prove that upper semicontinuity need not be imposed as a separate condition: it follows from subharmonicity along affine complex lines.

I designed the outline of this book and determined the material to be included. Most of the drafting was carried out by gpt-5.6-sol and Opus 5.5. I then revised and edited their drafts to produce the present text.

I am also working towards a complete formalization of the contents of this book in Lean 4. I hope to complete this work before the book is formally published, so that every proof in the published version is formally verified and no gaps remain.
